\scanpage{96}
\mspage{47}
donné $G$, l'ens $\widetilde A(G)$ et $I(G)$ (munis l'un et autre d'une application \inserted{``origine''} à valeurs dans $S(G)$).

Pour ceci, on a à utiliser la notion de ``graphes avec branches \unsure{orphelines|aplaties}'', qui généralise celle de graphe, et topologiquement correspond à la classification des \deleted{\unsure{4}} couples $(X,S)$ d'un espace topologique compact $X$, et d'une partie finie $S$, tels que a) $X$ soit une étoile au voisinage de tt pt de $S$, \deleted{et} b) \deleted{\unsure{S}} tt composante connexe de $X\setminus S$ est une $1$-variété à bord non compacte (car on veut que \unsure{toute|tout} \deleted{adhérence} d'arête ait un sommet). Il revient au $=$ de dire que $X$ est un $1$-complexe top, et $S$ une partie finie qui contient les sommets d'ordre $\geq3$, et telle qu'aucune composante connexe de $X\setminus S$ ne soit compacte (i.e. que tt composante connexe de $X$ contienne un sommet, i.e. $S\longrightarrow\pi_0(X)$ est surjectif).

Ceci se classifie combinatoirement par des systèmes $(S,\unsure{\widehat A|\widehat I},\sigma_{\unsure{\widehat A|\widehat I}},o_{\unsure{\widetilde A|\widehat A}})$, où \deleted{$\unsure{\widetilde A=\sigma_{\widetilde A}}$} $S,\unsure{\widehat A|\widehat I}$ sont des ens finis, $\sigma_{\widehat A}$ est une involution de $\widehat A$, et $o_{\unsure{\widetilde A|\widehat A}}:\widehat A\unsure{-|\longrightarrow}S$ une
\margin{95}
